% --- Title ---
\begin{frame}[plain]
\titlepage
\end{frame}

% --- Motivation ---
\begin{frame}{Motivation}
\begin{itemize}
  \item Industrial alloys (steels, Al alloys) are selected by their \textbf{elastic properties}: stiffness vs.\ elongation.
  \item Experimental measurement of elastic properties is \textbf{expensive and time-consuming}.
  \item Molecular simulations can predict elastic properties computationally---but they require accurate \textbf{interatomic potentials}.
  \item The Modified Embedded Atom Method (MEAM) is widely used for metals and alloys \cite{baskes1992, lee2001}.
\end{itemize}

\vspace{1em}
\begin{block}{Goal}
Build an automated pipeline to generate MEAM potentials for \emph{arbitrary} multi-element alloys.
\end{block}
\end{frame}

% --- The Problem ---
\begin{frame}{The Problem: Fragmented MEAM Libraries}
Published MEAM files cover only \textbf{subsets} of elements. No single file spans all elements needed for industrial alloys (e.g.\ Al\,7075: Al, Zn, Mg, Cu, Cr, Mn, Fe, Si, Ti).

\vspace{0.5em}
\centering
\small
\begin{tabular}{l|cccccccccccc}
\toprule
\textbf{Library File} & Al & Co & Cr & Cu & Fe & Mg & Mn & Mo & Ni & Si & Ti & Zn \\
\midrule
AlMgZn           & \checkmark &   &   &   &   & \checkmark &   &   &   &   &   & \checkmark \\
AlSiCuMgFe       & \checkmark &   &   & \checkmark & \checkmark & \checkmark &   &   &   & \checkmark &   &   \\
CuMo             &   &   &   & \checkmark &   &   &   & \checkmark &   &   &   &   \\
FeMnNiTiCuCrCoAl & \checkmark & \checkmark & \checkmark & \checkmark & \checkmark &   & \checkmark &   & \checkmark &   & \checkmark &   \\
\bottomrule
\end{tabular}

\vspace{1em}
Cross-interaction parameters between files are \textbf{undefined} $\Rightarrow$ cannot simulate combined compositions.
\end{frame}

% --- Research Question ---
\begin{frame}{Research Question}
\begin{block}{}
\centering
\emph{Can density functional theory reference calculations with neural-network surrogate optimization produce multi-element MEAM interatomic potentials that accurately predict the elastic properties of arbitrary metallic alloy compositions, even when the constituent elements do not share a common published MEAM parameterization?}
\end{block}
\end{frame}
